% sections/09-code-examples.tex — Code Examples

\section{Code Examples}

This section provides examples of expected documentation standards and code organization. These examples use unrelated classes to demonstrate the format only.

% ── Doxygen Documentation Format ──
\subsection{Doxygen Documentation Format}

\begin{lstlisting}[style=cpp]
/**
 * @file Dragon.h
 * @brief Declaration of the Dragon class.
 * @author Student Name
 * @date 2025-12-01
 */

/**
 * @brief Represents a mythical dragon creature.
 *
 * This class models a dragon with attributes such as fire power,
 * treasure hoard, and flight capability. It demonstrates proper
 * documentation format for class declarations.
 *
 * @tparam T The type of treasure the dragon hoards.
 */

/**
 * @brief Makes the dragon breathe fire at a target.
 *
 * @pre The dragon's fire power must be greater than zero.
 * @post The dragon's fire power is reduced by the attack cost.
 * @param target The name of the target to attack.
 * @return int The damage dealt to the target.
 * @throws std::runtime_error If the dragon has no fire power remaining.
 */
\end{lstlisting}

% ── Header File Structure ──
\subsection{Header File Structure}

\begin{lstlisting}[style=cpp, caption={Example Header: include/creatures/Dragon.h}]
/**
 * @file Dragon.h
 * @brief Declaration of the Dragon class.
 */

#pragma once

#include <string>
#include <memory>

namespace creatures {

/**
 * @brief Represents a fire-breathing dragon.
 *
 * A dragon can fly, breathe fire, and hoard treasure.
 * Each dragon has a unique name and a fire power level.
 */
class Dragon {
public:
    /**
     * @brief Constructs a new Dragon with the given name and fire power.
     *
     * @pre name must not be empty.
     * @post A Dragon object is created with the specified attributes.
     * @param name The unique name of the dragon.
     * @param firePower The initial fire power level.
     */
    explicit Dragon(const std::string& name, int firePower);

    /**
     * @brief Destroys the Dragon object.
     *
     * @post All resources owned by the dragon are released.
     */
    ~Dragon();

    // Rule of Five
    Dragon(const Dragon& other);
    Dragon& operator=(const Dragon& other);
    Dragon(Dragon&& other) noexcept;
    Dragon& operator=(Dragon&& other) noexcept;

    /**
     * @brief Makes the dragon breathe fire.
     *
     * @pre getFirePower() > 0.
     * @post Fire power is reduced by the attack cost.
     * @return int The damage dealt.
     * @throws std::runtime_error If fire power is zero or less.
     */
    int breatheFire();

    /**
     * @brief Adds treasure to the dragon's hoard.
     *
     * @pre amount > 0.
     * @post The dragon's treasure is increased by amount.
     * @param amount The amount of treasure to add.
     */
    void addTreasure(int amount);

    /**
     * @brief Allows the dragon to rest and recover fire power.
     *
     * @post Fire power is restored to its maximum value.
     */
    void rest();

    // Getters
    const std::string& getName() const noexcept;
    int getFirePower() const noexcept;
    int getTreasure() const noexcept;
    bool canFly() const noexcept;

private:
    struct Impl;                    ///< Forward-declared implementation.
    std::unique_ptr<Impl> pImpl_;   ///< Pointer to implementation.
};

} // namespace creatures
\end{lstlisting}

% ── Source File Structure ──
\subsection{Source File Structure}

\begin{lstlisting}[style=cpp, caption={Example Source: src/creatures/Dragon.cpp}]
/**
 * @file Dragon.cpp
 * @brief Implementation of the Dragon class.
 */

#include "creatures/Dragon.h"
#include <stdexcept>

namespace creatures {

struct Dragon::Impl {
    std::string name;
    int firePower;
    int treasure;
    bool flying;

    Impl(const std::string& n, int fp)
        : name(n), firePower(fp), treasure(0), flying(true) {}
};

Dragon::Dragon(const std::string& name, int firePower) {
    if (name.empty()) {
        throw std::invalid_argument("Dragon name cannot be empty");
    }
    pImpl_ = std::make_unique<Impl>(name, firePower);
}

Dragon::~Dragon() = default;

Dragon::Dragon(const Dragon& other)
    : pImpl_(std::make_unique<Impl>(*other.pImpl_)) {}

Dragon::Dragon(Dragon&& other) noexcept = default;

Dragon& Dragon::operator=(const Dragon& other) {
    if (this != &other) {
        pImpl_ = std::make_unique<Impl>(*other.pImpl_);
    }
    return *this;
}

Dragon& Dragon::operator=(Dragon&& other) noexcept = default;

int Dragon::breatheFire() {
    if (pImpl_->firePower <= 0) {
        throw std::runtime_error("No fire power remaining");
    }
    int damage = pImpl_->firePower * 2;
    pImpl_->firePower -= 10;
    return damage;
}

void Dragon::addTreasure(int amount) {
    if (amount > 0) {
        pImpl_->treasure += amount;
    }
}

void Dragon::rest() {
    pImpl_->firePower = 100;
}

const std::string& Dragon::getName() const noexcept {
    return pImpl_->name;
}

int Dragon::getFirePower() const noexcept {
    return pImpl_->firePower;
}

int Dragon::getTreasure() const noexcept {
    return pImpl_->treasure;
}

bool Dragon::canFly() const noexcept {
    return pImpl_->flying;
}

} // namespace creatures
\end{lstlisting}

% ── Main.cpp Structure ──
\subsection{Main.cpp Structure}

\begin{lstlisting}[style=cpp, caption={Example: src/main.cpp}]
/**
 * @file main.cpp
 * @brief Main entry point for the application.
 */

#include <iostream>
#include <string>

/**
 * @brief Prints usage information for the application.
 */
void printUsage() {
    std::cout << "Usage: ./app [options] <input_file>\n"
              << "Options:\n"
              << "  -h           Show this help message\n"
              << "  -v           Enable verbose output\n"
              << "  -o <file>    Specify output file\n";
}

int main(int argc, char* argv[]) {
    std::string inputFile = "data/input_default.json";
    bool verbose = false;

    // Parse command-line arguments
    for (int i = 1; i < argc; ++i) {
        std::string arg = argv[i];
        if (arg == "-h") {
            printUsage();
            return 0;
        } else if (arg == "-v") {
            verbose = true;
        } else {
            inputFile = arg;
        }
    }

    try {
        // 1. Load input data
        // 2. Initialize data structures
        // 3. Process data
        // 4. Output results

        if (verbose) {
            std::cout << "Processing: " << inputFile << "\n";
        }

        std::cout << "Application completed successfully.\n";
    } catch (const std::exception& e) {
        std::cerr << "Error: " << e.what() << "\n";
        return 1;
    }

    return 0;
}
\end{lstlisting}

% ── Pre/Post Condition Examples ──
\subsection{Pre/Post Condition Examples}

\begin{lstlisting}[style=cpp]
/**
 * @brief Feeds the dragon with the given food.
 *
 * @pre food must not be nullptr.
 * @pre food->isEdible() must return true.
 * @pre getHealth() < getMaxHealth().
 * @post getHealth() is increased by food->getNutrition().
 * @post getHealth() <= getMaxHealth().
 * @post food is consumed and ownership is transferred.
 * @param food Pointer to the food item (ownership transferred).
 * @throws std::invalid_argument If food is nullptr or not edible.
 * @throws std::logic_error If the dragon is already at full health.
 */
void feed(Food* food);

/**
 * @brief Removes a scale at the given position.
 *
 * @pre getScaleCount() > 0.
 * @pre position < getScaleCount().
 * @post getScaleCount() is decremented by one.
 * @post The removed scale is returned to the caller.
 * @param position The index of the scale to remove.
 * @return Scale* Pointer to the removed scale (caller owns memory).
 * @throws std::out_of_range If position >= getScaleCount().
 */
Scale* removeScale(size_t position);
\end{lstlisting}

\vspace{6pt}
\textbf{Important:} These examples demonstrate documentation format only. Your actual project must implement custom pointer-based data structures, not creature classes!
